\section{Referenciais e convenções adotadas}
\label{sec:Referenciais e convenções adotadas}

Foi utilizada a referência dextrogiro (\emph{LVLH, Local-Vertical-Local-Horizontal}) centrado no \emph{target} \cite{santosDextrogiro2021} para a formulação relativa e simulações numéricas.

\begin{figure}[ht] 
    \centering
    \includegraphics[width=0.5\textwidth]{2 Formulacao do problema/Imagens/LVLH.png}
    \caption{Referência LVLH centrado no target, adaptado de \cite{Folta2022}.}
    \label{fig:LVLH_NASA}
\end{figure}

A convenção
\footnote{Alguns autores, como \cite{fehse2003}, adotam a convenção alternativa na qual $\vec{x}$ é tangencial (\emph{along-track}) e $\vec{z}$ é radial.}
adotada neste trabalho segue a proposta de \cite{alfriend2010, Folta2022}, conforme:
o vetor $\vec{x}$ é orientado no sentido radial, apontando para fora da órbita a partir do centro de massa (CoM) da espaçonave, também indicando a altitude da espaçonave;
o vetor $\vec{y}$ é tangencial à órbita (também chamado de \emph{along-track}), apontando na direção do movimento orbital (vetor velocidade);
e o vetor $\vec{z}$ é normal ao plano orbital, alinhado ao vetor momento angular ($\vec \omega$); o nadir é a direção radial ($-R$).  

\subsection{Subsistema manipulador robótico}
O sistema de coordenadas do corpo (\emph{body frame}) é acoplado de forma rígida à plataforma onde o manipulador robótico está montado, portanto, não há movimento relativo entre eles. Esse sistema de coordenadas é referido como $\boldsymbol{\{B\}}$ — ou \emph{frame 0}, conforme a nomenclatura adotada por \cite{craig2005} — e possui sua origem coincidente com o centro de massa da plataforma. Dessa forma, o manipulador robótico compartilha a mesma translação e rotação da plataforma.

Essa referência será utilizada como base para a definição dos parâmetros inerciais da plataforma, tais como massa e matriz de inércia $\boldsymbol{I_B}$. Caso seja necessário utilizar em outra referência (por exemplo, fixado no sistema de coordenada inercial), uma transformação adicional deverá ser aplicada entre a referência global e o \emph{frame 0}.

A identificação dos elos segue a notação $\{i\}$, onde $i=1$ representa a primeira junta e $i=5$ a última, que acopla a ferramenta no último elo. Essa mesma numeração será utilizada para definir a tabela de Denavit-Hartenberg (DH) apresentada adiante.

Além disso, o efetuador final (\emph{end-effector}) do manipulador robótico é representado pelo $frame \{E\}$, cuja origem está localizada no final da ponta da ferramenta.

\subsubsection{Cinemática do manipulador robótico}
Cinemática é a área da mecânica que estuda o movimento dos corpos sem considerar as forças que assim o provocam \cite{craig2005}.
Dois espaços principais
\footnote{Conforme \cite{kucuk2006robot}, outras representações podem ser feitas, por exemplo, por vetores de Gibbs, parâmetros de Cayley-Klein e outros, no entanto, não serão abordados.}
são utilizados para modelar a cinemática de manipuladores robóticos: o espaço cartesiano e o espaço dos quaternions (apresentado na seção \eqref{sec:Quaternions}). Quando ocorre uma transformação entre dois sistemas de coordenadas cartesianos, essa transformação pode ser decomposta em uma rotação e uma translação \cite{kucuk2006robot}.  

Além disso, a cinemática pode ser classificada em \emph{cinemática direta} e \emph{cinemática inversa}. A cinemática direta normalmente é mais simples — envolvendo equações fechadas. Já a cinemática inversa tende a ser mais complexa, pois pode apresentar múltiplas soluções, ou até mesmo nenhuma solução válida. Essas soluções também podem envolver singularidades e não-linearidades, o que adiciona uma camada de complexidade ao problema \cite{craig2005, kucuk2006robot}.

\subsubsection{Cinemática direta do manipulador robótico}
Um manipulador robótico é composto por elos (\emph{links}) conectados em série por juntas (\emph{joints}) que permitem rotação ou translação relativa entre os elos.  

A cinemática direta consiste em determinar a posição e a orientação do efetuador final (\emph{end-effector}) em relação à referência da base (ou do corpo), a partir dos ângulos conhecidos das juntas \cite{kucuk2006robot}.

Para descrever essa cinemática, foram utilizados os parâmetros clássicos de Denavit-Hartenberg (DH), conforme originalmente propostos em \cite{denavitHartenberg1955}. Essa parametrização permite atribuir um sistema de coordenadas a cada elo. A abordagem clássica foi adotada neste trabalho pois não é comum encontrar em outras literaturas a sua forma original \cite{klug2018DH}.

O método de Denavit-Hartenberg utiliza quatro parâmetros para descrever, de forma sistemática e organizada, a posição e a orientação de cada elo do manipulador robótico \cite{denavitHartenberg1955, kucuk2006robot}. São eles:

\begin{itemize}
    \item \textbf{Ângulo da junta} ($\theta_i$): ângulo de rotação da junta $J_i$, em torno do eixo $z_{i-1}$;
    \item \textbf{Deslocamento axial} ($d_i$): distância ao longo do eixo $z_{i-1}$ entre as origens dos frames $i-1$ e $i$;
    \item \textbf{Comprimento do elo} ($a_i$): distância ao longo do eixo $x_i$ entre os eixos $z_{i-1}$ e $z_i$;
    \item \textbf{Ângulo de torção} ($\alpha_i$): ângulo entre os eixos $z_{i-1}$ e $z_i$, medido ao redor do eixo $x_i$.
\end{itemize}